\documentclass[12pt]{article}
\usepackage[utf8]{inputenc}
\usepackage[T1]{fontenc}
\usepackage[margin=1in]{geometry}
\usepackage{setspace}
\usepackage{xcolor}
\usepackage[colorlinks=true, urlcolor=blue!60!black, linkcolor=blue!60!black]{hyperref}
\setlength{\parindent}{0em}
\setlength{\parskip}{0.6em}
\renewcommand{\baselinestretch}{1.0}
\begin{document}
\begin{center}
{\LARGE \textit{Buckley v.\ Valeo}, 424 U.S.\ 1 (1976)\par}
\vspace{0.5em}
{\large Case Brief --- Two-Page Edition\par}
\end{center}

\section{Memory Jogger}
Contribution limits, disclosure, and voluntary public financing upheld; expenditure limits struck down; congressionally appointed FEC members invalid (\href{https://supreme.justia.com/cases/federal/us/424/1/}{Justia}).

\section{Detailed Case Facts}
After Watergate, the 1974 FECA amendments capped contributions (\$1,000 per candidate per election), capped independent, personal, and overall campaign spending, required disclosure of contributors over \$100, created public financing of presidential campaigns, and created an FEC most of whose voting members were appointed by congressional leaders. An ideologically diverse coalition led by Senator James Buckley and Eugene McCarthy sued.

\section{Procedural History}
The en banc D.C.\ Circuit upheld nearly all of FECA, treating money as regulable conduct under \textit{O'Brien}; a three-judge court upheld public financing. The Supreme Court affirmed in part and reversed in part.

\section{Judicial Votes}
Per curiam, with shifting majorities among eight Justices: contribution limits upheld 6--2 (Burger and Blackmun dissenting); expenditure limits struck 7--1 (White dissenting); disclosure upheld 7--1; public financing upheld; FEC appointments struck 7--1. Stevens took no part.

\section{Holding}
Spending money on politics is speech. Contribution limits are a marginal restriction valid if closely drawn to preventing \textit{quid pro quo} corruption or its appearance; expenditure limits fail exacting scrutiny. Disclosure and voluntary public financing are valid. Officers exercising significant federal authority must be appointed under Article~II.

\section{Analysis of Opinions}
\textbf{Per Curiam:} contributions are a ``general expression of support'' carrying corruption risk; independent expenditures pose a ``substantially diminished potential for abuse,'' and equalizing voices is no valid interest. \textbf{Burger:} ``contributions and expenditures are two sides of the same First Amendment coin.'' \textbf{White:} ``money is not speech.'' \textbf{Marshall:} the personal-funds ruling favors the rich. \textbf{Rehnquist:} public financing entrenched the major parties.

\section{Examples: Future Applications}
\textbf{Same-Side:} a state's \$1,075 contribution cap is valid under \textit{Nixon v.\ Shrink Missouri}. \textbf{Opposite-Side:} a minor party with documented harassment of supporters is exempt from donor disclosure. \textbf{Fence-Sitter:} public matching funds triggered by an opponent's spending, which add speech yet penalize expenditures (struck down in \textit{Arizona Free Enterprise}).

\section{Critique}
Critics say the contribution--expenditure line is unstable and pushed money into super PACs and ``dark money.'' Progressives fault the equation of money and speech; originalists, following Burger and Blackmun, see no principled difference between contributions and expenditures.

\section{Key Quotations}
\begin{quote}
``[T]he concept that government may restrict the speech of some elements of our society in order to enhance the relative voice of others is wholly foreign to the First Amendment.''
\end{quote}
\noindent --- Per Curiam, 424 U.S.\ at 48--49.

\end{document}
